% Generated draft: requires mathematical and source review.
\documentclass[11pt]{article}
\usepackage[margin=25mm]{geometry}
\usepackage{fontspec}
\setmainfont{Latin Modern Roman}
\usepackage{amsmath,amssymb,mathtools,unicode-math}
\setmathfont{Latin Modern Math}
\usepackage{enumitem,needspace,xurl,hyperref}
\setlength{\parindent}{0pt}
\setlength{\parskip}{5pt}
\setlength{\emergencystretch}{3em}
\newcommand{\N}{\mathbb{N}}
\newcommand{\Z}{\mathbb{Z}}
\newcommand{\Q}{\mathbb{Q}}
\newcommand{\R}{\mathbb{R}}
\newcommand{\C}{\mathbb{C}}
\newcommand{\F}{\mathbb{F}}
\newcommand{\A}{\mathbb{A}}
\newcommand{\PP}{\mathbb{P}}
\newcommand{\E}{\mathbb{E}}
\newcommand{\eps}{\varepsilon}
\newcommand{\dd}{\,\mathrm{d}}
\begin{document}
\section*{UnsolvedMath ID: 30003798}
\textbf{Source problem code:} catalogue entry
\subsection{Definitions and mathematical statement}
Let \(K\) be a field finitely generated over its prime subfield, and let \(G\) be an absolutely simple simply connected algebraic group over \(K\). Here absolutely simple means that \(G\) is connected and semisimple and, after extension to an algebraic closure of \(K\), has no nontrivial proper connected normal algebraic subgroup. Simply connected is understood in the algebraic-group sense: every central isogeny \(H\to G\) from a connected semisimple algebraic \(K\)-group is an isomorphism.

The supplied question is whether
\[
\text{for every such pair }(K,G),\qquad
K\text{ has good characteristic for }G
\ \Longrightarrow\ 
\bigl|\operatorname{gen}_{K}(G)\bigr|<\infty.
\]
The source record does not define either \(\operatorname{gen}_{K}(G)\) or the condition ``good characteristic.'' Consequently, this displayed formulation preserves the supplied question but does not yet specify a complete mathematical proposition.

One conventional interpretation of genus uses maximal tori. A maximal \(K\)-torus of \(G\) is a \(K\)-subgroup whose extension to an algebraic closure is a maximal torus. Under this interpretation, two groups have the same maximal \(K\)-tori if every maximal \(K\)-torus of either group is \(K\)-isomorphic to a maximal \(K\)-torus of the other. A possible definition is then
\[
\operatorname{gen}_{K}(G)=
\left\{
[G']_{K}:
\begin{array}{l}
G'\text{ is an inner }K\text{-form of }G,\\
G'\text{ and }G\text{ have the same maximal }K\text{-tori}
\end{array}
\right\},
\]
where \([G']_{K}\) denotes a \(K\)-isomorphism class, and an inner form is a form whose twisting cocycle takes values in the inner automorphism group of \(G\). This definition is an explicitly conditional interpretation, not a definition established by the supplied record. Both the intended genus convention and the exact characteristic restriction require source confirmation.

\textbf{Record qualification.} The supplied registry marks this statement as \texttt{unrecoverable} and describes UnsolvedMath ID \(30003798\), source code \texttt{catalogue entry}, as a duplicate extraction of a particular conjectural case retained in \texttt{OWR-16162-013}. It expressly declines to promote a separate clean statement. The text above therefore records the original question and its unresolved conventions; it does not establish a distinct conjecture or reconstruct the other record.
\subsection{Short English statement}
The recorded question asks whether every absolutely simple simply connected algebraic group over a finitely generated field of suitable characteristic has finite genus. The genus convention and the meaning of suitable characteristic are unspecified. The supplied registry treats this entry as a duplicate with an unrecoverable standalone statement; its open label is supplied curation, not an independently verified status.
\subsection{Sources}
\begin{itemize}
\item Local UnsolvedMath registry, permanent ID \(30003798\), source code \texttt{catalogue entry}, title \emph{Finiteness of Genera of Simply Connected Groups}, in the Oberwolfach Reports Open Problems collection. The supplied record gives proposed year \(2018\), names no individual proposer, and identifies this entry as a duplicate of a particular conjectural case retained in \texttt{OWR-16162-013}. Its supplied statement status is \texttt{unrecoverable}; the accompanying verification note is a registry assertion, not a verification performed for this draft.
\item Supplied report citation: \emph{Quadratic Forms and Related Structures over Fields} (2019), Oberwolfach Reports, DOI \texttt{10.4171/owr/2018/21}, pp.\ 1235--1291. \url{https://doi.org/10.4171/owr/2018/21}. The supplied background describes these reports as records of mathematical workshops and their open-problem discussions.
\item Supplied literature source: \emph{Groups with good reduction and genus}. \url{https://par.nsf.gov/servlets/purl/10588037}.
\item Supplied literature source: \emph{Oberwolfach Report 21}. \url{https://doi.org/10.4171/owr/2018/21}.
\item Literature assessment dated exactly \texttt{2026-08-21}: the supplied curation labels the problem open and says that finiteness of the genus for absolutely simple simply connected groups over finitely generated fields of good characteristic is still presented as a conjecture in current primary work. This is supplied dated curation derived from the OWR material, not an independently verified literature assessment. The supplied background explicitly requests independent verification of time-sensitive status and cited claims.
\end{itemize}
\end{document}
